% \section*{Глава 2. Проектирование веб-приложения}
% \addcontentsline{toc}{section}{Глава 2. Проектирование веб-приложения}
% \stepcounter{section}

% \subsection{Анализ целевой аудитории}
% Основными пользователями разрабатываемой системы являются покупатели спортивных
% товаров и администраторы интернет-магазина. Выделяются две ключевые роли:
% \begin{itemize}
% \item \textbf{Покупатель} --- зарегистрированный пользователь, осуществляющий
% поиск и выбор товаров, оформление заказов, а также управление личным
% кабинетом и просмотр истории покупок.
% \item \textbf{Администратор} --- сотрудник магазина, ответственный за
% управление каталогом товаров, обработку и отслеживание статусов
% заказов, а также работу с учётными записями пользователей.
% \end{itemize}
% Типичный покупатель --- человек в возрасте от 18 до 45 лет, активно занимающийся
% спортом или фитнесом, имеющий опыт онлайн-покупок и ожидающий интуитивно
% понятного интерфейса, быстрого поиска нужного товара и удобного оформления
% заказа. Ключевыми требованиями для данной роли являются: наглядный каталог с
% фильтрацией по категориям и характеристикам, корзина с возможностью изменения
% состава, а также личный кабинет с историей заказов и редактированием профиля.
% Администратор, как правило, является сотрудником компании со средним уровнем
% технической грамотности. Его основные задачи --- своевременное обновление
% ассортимента, контроль поступающих заказов и изменение их статусов по мере
% выполнения. Интерфейс административной панели должен обеспечивать быстрый
% доступ к ключевым операциям и минимизировать вероятность ошибок при управлении
% данными.
% Для каждой роли проведено профилирование по методу персонажей (Personas) для
% минимизации когнитивной нагрузки на интерфейс.


% \subsection{Описание функциональности приложения}
% Функциональные возможности системы определяются через взаимодействие пользователя
% с системой. Основной бизнес-кейс представлен в виде \textit{User Story}: «Я, как
% покупатель, хочу добавить товары в корзину и оформить заказ, чтобы получить
% выбранные спортивные товары с доставкой на указанный адрес».
% На рисунке~\ref{fig:use_case} представлена диаграмма вариантов использования
% (Use Case).
% \begin{figure}[ht!]
% 	\centering
% 	\includegraphics[width=0.8\textwidth]{use_case_diagram}
% 	\caption{Диаграмма вариантов использования системы}
% 	\label{fig:use_case}
% \end{figure}



% \subsection{Проектирование модели данных}
% При выборе архитектуры базы данных был проведён сравнительный анализ эффективности
% различных решений, опирающийся на данные о производительности современных
% СУБД~\cite{petrov2024}.
% Структура таблицы заказов представлена в табл.~\ref{tab:db_orders}.

% \begin{longtable}{|l|l|p{3.5cm}|p{5.5cm}|}
% 	\caption{Структура таблицы заказов (\texttt{orders})} \label{tab:db_orders}
% 	\hline
% 	\textbf{Поле} & \textbf{Тип} & \textbf{Ограничение} & \textbf{Описание} \ \hline
% 	\endfirsthead
% 	\tabcont{\ref{tab:db_orders}}
% 	\textbf{Поле} & \textbf{Тип} & \textbf{Ограничение} & \textbf{Описание} \ \hline
% 	\endhead
% 	\hline
% 	\endfoot
% 	\hline
% 	\endlastfoot
% 	id & UUID & PRIMARY KEY & Уникальный идентификатор заказа \ \hline
% 	user_id & UUID & FOREIGN KEY & Ссылка на таблицу пользователей \ \hline
% 	status & VARCHAR & NOT NULL & Статус (новый, в обработке, доставлен) \ \hline
% 	total_price & NUMERIC & NOT NULL & Итоговая стоимость заказа, руб. \ \hline
% 	address & TEXT & NOT NULL & Адрес доставки \ \hline
% 	contact_phone& VARCHAR & NOT NULL & Телефон получателя \ \hline
% 	comment & TEXT & NULL & Комментарий к заказу \ \hline
% 	created_at & TIMESTAMP & DEFAULT NOW() & Время оформления заказа \ \hline
% 	updated_at & TIMESTAMP & DEFAULT NOW() & Время последнего изменения \ \hline
% 	\end{longtable}
% 	Итоговая стоимость заказа рассчитывается следующим образом:
% 	\begin{equation}
% 	C_{total} = \sum_{i=1}^{n} (Q_i \cdot P_i) \cdot (1 - D) + S_{delivery},
% 	\label{eq:total_cost}
% 	\end{equation}
% 	\begin{where}
% 	\item $Q_i$~--- количество единиц $i$-го товара в заказе;
% 	\item $P_i$~--- цена единицы $i$-го товара, руб.;
% 	\item $D$~--- коэффициент скидки (от 0 до 1);
% 	\item $S_{delivery}$~--- стоимость доставки, руб.
% \end{where}
% На рисунке~\ref{fig:er_diagram} представлена ER-диаграмма базы данных,
% отражающая сущности системы и связи между ними.
% \begin{figure}[ht!]
% \centering
% \includegraphics[width=0.95\textwidth]{er_diagram}
% \caption{ER-диаграмма базы данных интернет-магазина}
% \label{fig:er_diagram}
% \end{figure}
% Логическая схема базы данных детализирует атрибуты каждой сущности и типы
% связей между ними (один-ко-многим, многие-ко-многим), не привязываясь к
% конкретной СУБД. Физическая схема, представленная на
% рисунке~\ref{fig:physical_schema}, специфицирует конкретные типы данных,
% индексы и внешние ключи применительно к выбранной СУБД PostgreSQL.
% \begin{figure}[ht!]
% 	\centering
% 	\includegraphics[width=0.95\textwidth]{physical_schema}
% 	\caption{Физическая схема базы данных}
% 	\label{fig:physical_schema}
% \end{figure}



% \subsection{Разработка макетов страниц}

% *Макеты страниц*